\documentclass[10pt,a4paper]{amsart}
\usepackage[margin=19mm]{geometry}
\usepackage{amsmath,amssymb,mathtools}
\usepackage[scr=rsfs,cal=euler]{mathalfa}
\usepackage{xfrac}
\usepackage[unicode,hidelinks,bookmarks=false]{hyperref}
\newcommand{\ie}{i.e.\ }
\newcommand{\cf}{cf.\ }
\newcommand{\resp}{respectively\ }
\newcommand{\Subsection}{\subsection}
\providecommand{\nobreakdash}{\nobreak\hbox{-}}
\setlength{\emergencystretch}{2em}
\multlinegap0pt
\usepackage{bm}
\usepackage{bbm}
\usepackage{dsfont}
\usepackage{xspace}
\usepackage{cancel}
\usepackage{mathtools}
\usepackage{amsthm}
\makeatletter
\@ifundefined{theorem}{\newtheorem{theorem}{Theorem}}{}
\@ifundefined{lemma}{\newtheorem{lemma}[theorem]{Lemma}}{}
\makeatother
\DeclareMathOperator{\tr}{tr}
\newcommand{\R}{\mathbb{R}}
\newcommand{\bS}{\mathbf{S}}
\title[Accelerating operator Sinkhorn iteration with overrelaxation]{Accelerating operator Sinkhorn iteration with overrelaxation}
\author{Tasuku Soma, Andr\'e Uschmajew}
\date{Source: arXiv:2410.14104v2 (CC BY 4.0)}
\begin{document}
\maketitle

\noindent\emph{Source and licence.} Accelerating operator Sinkhorn iteration with overrelaxation. Version arXiv:2410.14104v2 (CC BY 4.0), distributed under the Creative Commons Attribution 4.0 International licence, as recorded in the arXiv licence metadata for this exact version and shown on the abstract page https://arxiv.org/abs/2410.14104v2. Full rights notice: licenses/arxiv-2410.14104-CC-BY-4.0.txt.\par\smallskip

\noindent\textit{Editorial scope.} Counted unit: the Lemma whose source label is \texttt{lem: null space of Hessian} in the article (source0.tex lines 410--412, source label \texttt{lem: null space of Hessian}), delivered together with the article's own complete original proof of it (source0.tex lines 414--448, one \texttt{proof} environment of 35 lines). No mathematics is written, repaired or re-derived: statement and proof are the article's own text line for line. Required context retained uncounted: eq: fixed-point map (lines 303--305). Every definition, display and label of those lines is retained unchanged. Omitted from this package and not redistributed: 1--302, 306--409, 413--413, 449--1140. Nothing inside a retained excerpt is omitted or altered: every retained body is delivered exactly as the article's lines are written, so no omission receipt of any kind accompanies this package. Citations: the retained excerpts carry no citation command at all, so no bibliography entry of the article is transcribed and no reference is redistributed. Reference adaptation: none was needed and none was made; every reference inside the delivered unit resolves inside it, so no unresolved reference or citation is produced. Printed slips kept literal: no printed word, symbol, hypothesis or number of the counted statement or of its proof was altered. Layout and class: the article's own class is \texttt{article} with packages including fullpage, fontenc, graphicx, amsmath, amssymb, amsthm, mathtools, microtype, xcolor, hyperref, algorithm, algorithmicx, algpseudocode, subcaption; the wrapper is the lane's pinned \texttt{amsart} editorial wrapper at 10pt on A4, which restates the theorem-like environments the retained text needs and adds the article's own macro definitions that the retained text needs (\texttt{\textbackslash R}, \texttt{\textbackslash bS}, \texttt{\textbackslash tr}), with extra wrapper packages (bm, bbm, dsfont, xspace, cancel, mathtools). The delivered numbering is the unit's own. Assets: no retained span contains a figure environment, a graphics-inclusion command, a PDF-page inclusion, a table or a TikZ picture; no image file is copied into this package, the package declares no asset dependency and no image file is redistributed with it. Licence: version arXiv:2410.14104v2 of this article is distributed under the Creative Commons Attribution 4.0 International licence, as recorded in the arXiv licence metadata for this exact version and shown on the abstract page https://arxiv.org/abs/2410.14104v2. A full rights notice is delivered at licenses/arxiv-2410.14104-CC-BY-4.0.txt. Characters: every retained excerpt is pure ASCII, so the delivered engine, XeLaTeX with its default Latin Modern text font, reports no missing character.
\begin{equation}\label{eq: fixed-point map}
\bS\!\begin{pmatrix} X \\ Y \end{pmatrix} = \begin{pmatrix} S_1(Y) \\ S_2(S_1(Y)) \end{pmatrix}.
\end{equation}

\begin{lemma}\label{lem: null space of Hessian}
Assume that the CP map $\Phi$ is positivity improving, and let $(X_*,Y_*)$ be a fixed-point of the map $\bS$ in~\eqref{eq: fixed-point map}. Then the Hessian $\nabla^2 f(X_*,Y_*)$ is positive semidefinite with a one-dimensional null space spanned by $(X_*,-Y_*)$. In other words, the null space of $\nabla^2 f(X_*,Y_*)$ equals the tangent space to the orbit $(\mu X_*,\mu^{-1} Y_*)$ at $\mu = 1$.
\end{lemma}

\begin{proof}
According to the considerations preceding the lemma, $(X_*,Y_*)$ is a global minimum of $f$. Therefore, the Hessian is positive semidefinite. Since $\nabla f$ is constant zero on the orbit $(\mu X_*, \mu^{-1} Y_*)$, it also follows that the tangent space to the orbit at $(X_*,Y_*)$ is in the null space of the Hessian. This can also be verified directly: as a linear map on $\R^{m \times m}_{\mathrm{sym}} \times \R^{n \times n}_{\mathrm{sym}}$ the Hessian at $(X_*,Y_*)$ reads
\[
\begin{pmatrix}
 H_1 \\ H_2
\end{pmatrix}
\mapsto
\begin{pmatrix}
 \Phi(H_2) + \frac{1}{m} X_*^{-1} H_1 X_*^{-1} \\ \Phi^*(H_1) + \frac{1}{n} Y_*^{-1} H_2 Y_*^{-1}
\end{pmatrix}.
\]
Using the fixed-point property $\Phi(Y_*) = \frac{1}{m} X_*^{-1}$ and $\Phi^*(X_*) = \frac{1}{n} Y_*^{-1}$ we see that $(H_1,H_2) = (X_*,-Y_*)$ is indeed in the null space of the Hessian. We will show that this is the only direction in the null space.

As can be derived from the above explicit expression of the Hessian, a necessary condition for $(H_1,H_2)$ to be in the null space is that the matrix $\tilde H_1 = X_*^{-1/2} H_1 X_*^{-1/2}$ satisfies the equation
\[
 \tilde H_1 = mn X_*^{1/2} \Phi \big( Y_* \Phi^*(X_*^{1/2} \tilde H_1 X_*^{1/2}) Y_* \big) X_*^{1/2} \eqqcolon \Gamma( \tilde H_1 ),
\]
that is, $\tilde H_1$ must be an eigenvector of the linear operator $\Gamma$ on $\R^{m \times m}_{\mathrm{sym}}$ with eigenvalue one. We know that $\tilde H_1 = I_m$ (corresponding to $H_1 = X_*$) is such an eigenvector, and the lemma will be proved if we show that (up to scalar multiplication) it is the only one. To do this, let
\begin{equation}\label{eq:expansion H_1}
 \tilde H_1 = \sum_{i=1}^m z_i^{} u_i^{} u_i^\top
\end{equation}
be a spectral decomposition with pairwise orthonormal vectors $u_1,\dots,u_m \in \R^m$ and scalars $z_1,\dots,z_m$, so that
\[
 \Gamma(\tilde H_1) = \sum_{i=1}^m z_i^{} \Gamma(u_i^{} u_i^\top).
\]
If $\tilde H_1 = \Gamma(\tilde H_1)$, then by comparing coefficients for the $u_1^{} u_1^\top,\dots,u_m^{} u_m^\top$ we obtain a relation
\[
 z = B z
\]
where the matrix $B = [b_{ij}]$ has entries
\[
 b_{ij} = \tr(\Gamma(u_j^{} u_j^\top) u_i^{} u_i^\top) = u_i^\top \Gamma(u_j^{} u_j^\top) u_i.
\]
Although this matrix depends on the orthonormal basis $u_1,\dots,u_m$, the vector $z = \mathbf 1$ always satisfies the above relation, that is, $B \mathbf 1 = \mathbf 1$, since $I_m = \sum_{i=1}^m u_i^{} u_i^\top$ for any choice of basis. The main observation now is that since $\Phi^*$ (together with $\Phi$) is positivity improving, $\Gamma(u_j^{} u_j^\top)$ is a positive definite matrix for every~$j$, implying that all entries of $B$ are strictly positive. We can then deduce from the Perron--Frobenius theorem that $\mathbf 1$ is the only eigenvector of $B$ with eigenvalue one. This shows that for any orthonormal basis $u_1,\dots,u_m$, $\tilde H_1 = I_m$ is the only eigenvector of $\Gamma$ of the form~\eqref{eq:expansion H_1} with eigenvalue one, which concludes the proof.
\end{proof}

\end{document}
